% ELAW 技术栈调研报告
% 使用XeLaTeX + ctexart编译
\documentclass[fontset=mac,a4paper,11pt]{ctexart}

% ===== 基础宏包 =====
\usepackage{geometry}
\geometry{top=2.5cm,bottom=2.5cm,left=2.5cm,right=2.5cm}
\usepackage{graphicx}
\usepackage{xcolor}
\usepackage{booktabs}
\usepackage{longtable}
\usepackage{enumitem}
\usepackage{hyperref}
\usepackage{fancyhdr}
\usepackage{titlesec}
\usepackage{listings}
\usepackage{tikz}
\usepackage{float}
\usepackage{amsmath}
\usepackage{amssymb}
\usepackage{makecell}
\usepackage{multirow}
\usepackage{colortbl}
\usepackage{array}

% TikZ库
\usetikzlibrary{shapes.geometric, arrows.meta, positioning, fit, backgrounds, calc, shadows}

% ===== 颜色定义 =====
\definecolor{primaryblue}{RGB}{0,102,204}
\definecolor{darkblue}{RGB}{0,51,102}
\definecolor{lightbg}{RGB}{245,248,252}
\definecolor{codebg}{RGB}{248,248,248}
\definecolor{codeframe}{RGB}{200,200,200}
\definecolor{risklow}{RGB}{144,190,109}
\definecolor{riskmedium}{RGB}{249,166,62}
\definecolor{riskhigh}{RGB}{217,83,79}
\definecolor{tableheader}{RGB}{240,242,245}
\definecolor{accentgreen}{RGB}{34,139,34}

% ===== Hyperref设置 =====
\hypersetup{
    colorlinks=true,
    linkcolor=primaryblue,
    urlcolor=primaryblue,
    citecolor=primaryblue,
    pdfauthor={ELAW Team},
    pdftitle={ELAW 技术栈调研报告}
}

% ===== 代码高亮设置 =====
\lstset{
    basicstyle=\ttfamily\small,
    backgroundcolor=\color{codebg},
    frame=single,
    rulecolor=\color{codeframe},
    breaklines=true,
    showstringspaces=false,
    keywordstyle=\color{primaryblue}\bfseries,
    commentstyle=\color{gray}\itshape,
    stringstyle=\color{accentgreen},
    numbers=left,
    numberstyle=\tiny\color{gray},
    numbersep=8pt,
    xleftmargin=15pt,
    xrightmargin=5pt,
    aboveskip=10pt,
    belowskip=10pt,
    tabsize=2,
    captionpos=b
}

% ===== 页眉页脚 =====
\pagestyle{fancy}
\fancyhf{}
\fancyhead[L]{\small\leftmark}
\fancyhead[R]{\small ELAW 技术栈调研报告}
\fancyfoot[C]{\thepage}
\renewcommand{\headrulewidth}{0.4pt}
\renewcommand{\footrulewidth}{0pt}

% ===== 标题格式 =====
\titleformat{\section}
{\normalfont\Large\bfseries\color{darkblue}}
{\thesection}{1em}{}
\titleformat{\subsection}
{\normalfont\large\bfseries\color{primaryblue}}
{\thesubsection}{1em}{}
\titleformat{\subsubsection}
{\normalfont\normalsize\bfseries}
{\thesubsubsection}{1em}{}

% ===== 自定义表格样式 =====
\newcolumntype{L}[1]{>{\raggedright\arraybackslash}p{#1}}
\newcolumntype{C}[1]{>{\centering\arraybackslash}p{#1}}
\newcolumntype{R}[1]{>{\raggedleft\arraybackslash}p{#1}}

% ===== 风险矩阵单元格颜色 =====
\newcommand{\riskcell}[1]{%
\ifnum#1=1\cellcolor{risklow!30}低%
\else\ifnum#1=2\cellcolor{riskmedium!30}中%
\else\cellcolor{riskhigh!30}高%
\fi\fi
}

\begin{document}

% ===== 封面 =====
\begin{titlepage}
\begin{center}
\vspace*{2cm}
{\Huge\bfseries\color{darkblue} ELAW 技术栈调研报告}\\[0.8cm]
{\Large\color{gray} Exusiai Live Agent WebUI}\\[0.3cm]
{\large 从 ACP-QQ-Bridge 到浏览器架构的技术演进}\\[2.5cm]
\begin{tikzpicture}
\node[draw=primaryblue, line width=1.5pt, rounded corners=10pt, inner sep=20pt, fill=lightbg] {
\begin{minipage}{0.6\textwidth}
\centering
\textbf{项目代号：}ELAW\\[0.3cm]
\textbf{报告类型：}技术调研与架构设计\\[0.3cm]
\textbf{版本：}v1.0\\[0.3cm]
\textbf{日期：}\today
\end{minipage}
};
\end{tikzpicture}
\vfill
{\small\color{gray} 本报告由技术调研团队编写，基于现有项目文档与行业最佳实践}
\end{center}
\end{titlepage}

% ===== 目录 =====
\tableofcontents
\newpage

% ============================================================
\section{项目技术需求概述}
% ============================================================

\subsection{项目背景与演进}

ELAW (Exusiai Live Agent WebUI) 是 ACP-QQ-Bridge 项目的完全重构。旧项目通过QQ机器人协议桥接Kimi Code CLI，实现了一个以QQ消息为交互界面的AI助手。新项目的核心转变是将交互界面从QQ迁移到浏览器，从而解锁了更丰富的UI能力和更直接的实时通信。

\begin{table}[H]
\centering
\caption{旧项目与新技术架构对比}
\label{tab:arch-evolution}
\begin{tabular}{L{3.5cm}L{5cm}L{5cm}}
\toprule
\textbf{维度} & \textbf{旧项目 (ACP-QQ-Bridge)} & \textbf{新项目 (ELAW)} \\
\midrule
\rowcolor{tableheader}
交互界面 & QQ消息气泡（文本+图片） & 浏览器响应式Web UI \\
通信协议 & QQ Bot API (HTTP轮询) & WebSocket全双工通信 \\
\rowcolor{tableheader}
实时性 & 消息轮询，延迟2-5秒 & 实时流式推送，延迟<100ms \\
多媒体支持 & 静态表情包、图片 & Live2D模型、TTS语音、动画 \\
\rowcolor{tableheader}
部署方式 & 云服务器+QQ账号 & 局域网PWA+浏览器 \\
用户场景 & 多用户QQ群聊 & 个人/本地单用户交互 \\
\bottomrule
\end{tabular}
\end{table}

\subsection{核心需求驱动因素}

新项目的技术选型由以下核心需求驱动：

\begin{enumerate}[leftmargin=2em]
\item \textbf{实时流式反馈}：用户输入后，AI的响应需要逐字/逐块显示在单条消息气泡内，而非等待完整响应后才展示。这要求WebSocket长连接 + 前端增量渲染。

\item \textbf{Live2D角色交互}：需要加载和驱动Live2D模型，实现表情切换、动作播放、口型同步（lip sync）。这要求WebGL支持、音频分析能力。

\item \textbf{打断机制}：用户可以在AI生成回复的过程中发送新消息，中断当前生成。这要求前后端都支持取消/重置状态。

\item \textbf{响应式多端支持}：桌面端和手机端都需良好体验，且手机端应支持PWA安装到主屏幕。

\item \textbf{本地化优先}：阶段1完全在局域网内运行，无需公网服务器，减少外部依赖。
\end{enumerate}

% ============================================================
\section{前端技术栈详解}
% ============================================================

\subsection{框架选型：React 18 + TypeScript}

\subsubsection{候选框架对比}

\begin{table}[H]
\centering
\caption{前端框架对比分析}
\label{tab:frontend-framework}
\begin{tabular}{L{2.5cm}L{3.5cm}L{3.5cm}L{3.5cm}}
\toprule
\textbf{维度} & \textbf{React 18} & \textbf{Vue 3} & \textbf{Svelte} \\
\midrule
\rowcolor{tableheader}
生态系统 & 极其丰富，第三方库覆盖全面 & 丰富，中文社区活跃 & 较新，库数量较少 \\
Live2D集成 & 成熟库（React绑定完善） & 有但较分散 & 几乎无现成方案 \\
\rowcolor{tableheader}
TypeScript & 原生支持最好 & 支持良好 & 支持但有局限 \\
学习曲线 & 中等（需理解Hooks） & 较低（模板直观） & 低（接近原生JS） \\
\rowcolor{tableheader}
流式渲染 & Concurrent Features优秀 & 支持但不如React & 有限支持 \\
社区规模 & 最大（GitHub Stars 220k+） & 大（GitHub Stars 38k+） & 中等（75k+） \\
\rowcolor{tableheader}
招聘/人才 & 最容易 & 容易 & 较难 \\
\bottomrule
\end{tabular}
\end{table}

\subsubsection{选择React 18的理由}

\begin{itemize}[leftmargin=2em]
\item \textbf{Concurrent Features}：React 18的并发渲染特性对流式UI至关重要。AI回复逐字追加时，可以使用\texttt{useTransition}或\texttt{useDeferredValue}保证UI不因高频更新而卡顿。
\item \textbf{成熟的Live2D生态}：Live2D Cubism Web SDK已有成熟的React封装示例，且React的声明式组件模型便于管理Live2D模型的生命周期。
\item \textbf{TypeScript原生支持}：React的类型定义最完善，配合泛型组件可构建类型安全的UI系统。
\item \textbf{团队经验}：旧项目CLI桥接代码已有TypeScript经验，迁移成本低。
\end{itemize}

\subsection{构建与开发工具：Vite}

\begin{table}[H]
\centering
\caption{构建工具对比}
\label{tab:build-tools}
\begin{tabular}{L{2.5cm}L{3.5cm}L{3.5cm}L{3.5cm}}
\toprule
\textbf{维度} & \textbf{Vite} & \textbf{Webpack 5} & \textbf{Create React App} \\
\midrule
\rowcolor{tableheader}
冷启动速度 & 毫秒级（原生ESM） & 秒级（需打包） & 较慢（已弃用） \\
HMR速度 & 毫秒级 & 秒级 & 秒级 \\
\rowcolor{tableheader}
PWA支持 & vite-plugin-pwa原生支持 & 需配置workbox-webpack-plugin & 需手动配置 \\
配置复杂度 & 极简，开箱即用 & 复杂，需大量配置 & 不灵活，隐藏配置 \\
\rowcolor{tableheader}
TypeScript & 原生支持（esbuild） & 需配置ts-loader & 内置但过时 \\
生产构建 & Rollup，tree-shaking优秀 & 优化成熟但慢 & 已停止维护 \\
\rowcolor{tableheader}
社区趋势 & 主流（新项目首选） & 存量项目为主 & 官方已弃用 \\
\bottomrule
\end{tabular}
\end{table}

\textbf{Vite的关键优势：}

\begin{enumerate}[leftmargin=2em]
\item \textbf{开发速度}：基于原生ESM的冷启动和热更新，在AI助手这种需要频繁迭代UI的项目中效率极高。
\item \textbf{PWA一键支持}：\texttt{vite-plugin-pwa}自动生成Service Worker、manifest.json，并支持Workbox离线缓存策略。
\item \textbf{构建产物优化}：自动代码分割、CSS代码抽取、静态资源优化，生成的包体积小。
\end{enumerate}

\subsection{状态管理：Zustand}

\begin{table}[H]
\centering
\caption{状态管理方案对比}
\label{tab:state-management}
\begin{tabular}{L{2.5cm}L{3.5cm}L{3.5cm}L{3.5cm}}
\toprule
\textbf{维度} & \textbf{Zustand} & \textbf{Redux Toolkit} & \textbf{Context API} \\
\midrule
\rowcolor{tableheader}
包大小 & ~1KB & ~11KB+（RTK+React绑定） & 0（内置） \\
学习成本 & 极低 & 中等（需理解Slice等概念） & 低 \\
\rowcolor{tableheader}
TypeScript & 类型推导完美 & 需要写模板代码 & 需手动定义类型 \\
中间件支持 & 丰富（persist、devtools等） & 丰富 & 无 \\
\rowcolor{tableheader}
适用规模 & 中小项目（ELAW规模） & 大型复杂应用 & 简单全局状态 \\
性能 & 无不必要的重渲染 & 需要优化selector & 易触发不必要重渲染 \\
\rowcolor{tableheader}
异步逻辑 & 简洁直接 & 需要Thunk/Saga & 需在外层处理 \\
\bottomrule
\end{tabular}
\end{table}

Zustand适合本项目的原因：
\begin{itemize}[leftmargin=2em]
\item ELAW是单用户应用，状态规模可控，无需Redux的严格架构。
\item Zustand的API极简（类似React hooks），不需要Provider包裹，降低了心智负担。
\item 内置\texttt{persist}中间件可自动将配置状态持久化到localStorage。
\end{itemize}

\subsection{UI与样式方案：Tailwind CSS + CSS Variables}

\begin{table}[H]
\centering
\caption{样式方案对比}
\label{tab:styling}
\begin{tabular}{L{2.5cm}L{3.5cm}L{3.5cm}L{3.5cm}}
\toprule
\textbf{维度} & \textbf{Tailwind CSS} & \textbf{CSS-in-JS (styled-components)} & \textbf{SCSS/SASS} \\
\midrule
\rowcolor{tableheader}
开发速度 & 极快（原子类直接写） & 中等 & 中等（需定义类名） \\
包大小 & 生产时Purge后极小 & 运行时开销较大 & 纯CSS，无运行时 \\
\rowcolor{tableheader}
主题切换 & 与CSS Variables天然配合 & 需要ThemeProvider & 需要预定义变量 \\
学习成本 & 需记忆类名 & 低（标准CSS） & 低（标准CSS） \\
\rowcolor{tableheader}
Tailwind支持 & 自身 & 无（需额外配置） & 无 \\
\bottomrule
\end{tabular}
\end{table}

\textbf{选型方案：}以Tailwind CSS为主力，配合CSS Variables实现动态主题（如日间/夜间模式切换、角色主色调切换）。对于Live2D容器的特殊样式（如固定定位、半透明覆盖层），使用纯CSS模块作为补充。

\subsection{Markdown渲染与代码高亮：shiki}

\begin{table}[H]
\centering
\caption{代码高亮方案对比}
\label{tab:code-highlight}
\begin{tabular}{L{2.5cm}L{3.5cm}L{3.5cm}L{3.5cm}}
\toprule
\textbf{维度} & \textbf{shiki} & \textbf{prismjs} & \textbf{highlight.js} \\
\midrule
\rowcolor{tableheader}
渲染质量 & 极高（基于TextMate语法，VS Code同款） & 中等（正则匹配） & 中等（正则匹配） \\
主题支持 & 内置VS Code全主题 & 数量多但质量参差 & 数量多但质量参差 \\
\rowcolor{tableheader}
渲染方式 & 服务端/构建时渲染（SSR友好） & 客户端运行时 & 客户端运行时 \\
包大小 & 按需加载语言（极小） & ~4KB（核心）+ 语言包 & ~6KB（核心）+ 语言包 \\
\rowcolor{tableheader}
React集成 & react-shiki支持良好 & react-syntax-highlighter成熟 & 可用但非主流 \\
\bottomrule
\end{tabular}
\end{table}

shiki在本项目中的优势尤为明显：AI助手的代码展示需要与Kimi Code CLI（VS Code基）风格一致，shiki的TextMate语法高亮能够完美复现VS Code的视觉效果。

\subsection{Live2D集成技术：Cubism Web SDK}

\subsubsection{SDK架构概述}

Live2D Cubism Web SDK 由以下核心模块组成：

\begin{itemize}[leftmargin=2em]
\item \textbf{Core}：模型解析、参数管理、动作系统
\item \textbf{Framework}：渲染器抽象、纹理管理
\item \textbf{WebGL Renderer}：实际的WebGL绘制实现
\item \textbf{Motion Management}：动作文件（.motion3.json）加载与播放
\item \textbf{Expression}：表情文件（.exp3.json）动态插值
\item \textbf{LipSync}：口型同步（基于音频输入值映射到 mouth open 参数）
\end{itemize}

\subsubsection{移动端WebGL兼容性考量}

\begin{table}[H]
\centering
\caption{移动端WebGL兼容性分析}
\label{tab:webgl-mobile}
\begin{tabular}{L{3.5cm}L{4cm}L{4.5cm}}
\toprule
\textbf{平台/浏览器} & \textbf{WebGL 2.0支持} & \textbf{Live2D注意事项} \\
\midrule
\rowcolor{tableheader}
iOS Safari (16+) & 支持 & 纹理尺寸限制4096x4096 \\
Chrome Android & 支持 & 性能取决于GPU型号 \\
\rowcolor{tableheader}
MIUI浏览器 & 部分支持 & 可能需降级到WebGL 1.0 \\
微信内置浏览器 & 支持（iOS WKWebView） & 同Safari限制 \\
\rowcolor{tableheader}
iOS 15以下 & 仅WebGL 1.0 & Cubism SDK 4.x支持WebGL 1.0 \\
\bottomrule
\end{tabular}
\end{table}

\textbf{缓解策略：}实现WebGL 2.0/1.0自动降级检测；为低端设备提供静态PNG表情回退方案；模型纹理控制在2048x2048以内。

\subsection{PWA技术实现}

Vite PWA Plugin (\texttt{vite-plugin-pwa}) 提供了完整的PWA构建支持：

\begin{lstlisting}[caption=Vite PWA配置示例]
// vite.config.ts
import { VitePWA } from 'vite-plugin-pwa'

export default {
  plugins: [
    VitePWA({
      registerType: 'autoUpdate',
      manifest: {
        name: 'ELAW - Exusiai Agent',
        short_name: 'ELAW',
        start_url: '/',
        display: 'standalone',
        background_color: '#ffffff',
        theme_color: '#0066cc',
        icons: [
          { src: '/icon-192.png', sizes: '192x192' },
          { src: '/icon-512.png', sizes: '512x512' }
        ]
      },
      workbox: {
        globPatterns: ['**/*.{js,css,html,ico,png,svg,woff2}'],
        runtimeCaching: [
          {
            urlPattern: /^ws:\/\/.*/i,
            handler: 'NetworkOnly'
          }
        ]
      }
    })
  ]
}
\end{lstlisting}

\textbf{PWA关键配置说明：}
\begin{itemize}[leftmargin=2em]
\item \texttt{registerType: 'autoUpdate'}：Service Worker在检测到新版本时自动更新。
\item \texttt{display: 'standalone'}：添加到主屏幕后无浏览器地址栏，体验接近原生App。
\item WebSocket连接需配置为\texttt{NetworkOnly}，避免缓存问题。
\item 需要HTTPS（或localhost）才能注册Service Worker，局域网IP访问时Safari支持有限，Chrome完全支持。
\end{itemize}

\subsection{WebSocket客户端设计}

\begin{table}[H]
\centering
\caption{WebSocket客户端方案对比}
\label{tab:websocket-client}
\begin{tabular}{L{2.5cm}L{3.5cm}L{3.5cm}L{3.5cm}}
\toprule
\textbf{维度} & \textbf{原生 WebSocket} & \textbf{Socket.io Client} & \textbf{Reconnecting-WS} \\
\midrule
\rowcolor{tableheader}
包大小 & 0（内置） & ~30KB & ~3KB \\
重连机制 & 手动实现 & 内置（指数退避） & 内置（简单退避） \\
\rowcolor{tableheader}
心跳检测 & 手动实现 & 内置 & 手动实现 \\
二进制支持 & 原生支持 & 支持（封装） & 原生支持 \\
\rowcolor{tableheader}
服务端兼容性 & 通用 & 需Socket.io服务端 & 通用 \\
\bottomrule
\end{tabular}
\end{table}

\textbf{推荐方案：}原生WebSocket API + 自定义重连逻辑。理由：后端使用FastAPI原生WebSocket而非Socket.io，保持协议一致；原生API可精确控制重连策略（适合移动端场景）。

\begin{lstlisting}[caption=WebSocket客户端封装示例]
class ELAWWebSocketClient {
  private ws: WebSocket | null = null
  private reconnectAttempts = 0
  private maxReconnectAttempts = 10
  private heartbeatInterval: number | null = null
  private url: string

  constructor(url: string) {
    this.url = url
  }

  connect() {
    this.ws = new WebSocket(this.url)
    this.ws.onopen = () => {
      this.reconnectAttempts = 0
      this.startHeartbeat()
    }
    this.ws.onclose = () => {
      this.stopHeartbeat()
      this.scheduleReconnect()
    }
    this.ws.onmessage = (e) => this.handleMessage(e.data)
  }

  private scheduleReconnect() {
    if (this.reconnectAttempts >= this.maxReconnectAttempts) return
    const delay = Math.min(1000 * 2 ** this.reconnectAttempts, 30000)
    setTimeout(() => { this.reconnectAttempts++; this.connect() }, delay)
  }

  private startHeartbeat() {
    this.heartbeatInterval = window.setInterval(() => {
      this.send(JSON.stringify({ type: 'ping' }))
    }, 30000)
  }
  // ...
}
\end{lstlisting}

\subsection{动画与交互：Framer Motion}

\begin{table}[H]
\centering
\caption{动画方案对比}
\label{tab:animation}
\begin{tabular}{L{2.5cm}L{3.5cm}L{3.5cm}L{3.5cm}}
\toprule
\textbf{维度} & \textbf{Framer Motion} & \textbf{CSS Transition} & \textbf{GSAP} \\
\midrule
\rowcolor{tableheader}
React集成 & 原生组件（motion.div） & 需手动管理className & 命令式API，需封装 \\
声明式API & 完美支持 & 有限 & 无 \\
\rowcolor{tableheader}
手势支持 & 内置（drag、hover、tap） & 无 & 需额外插件 \\
AnimatePresence & 处理组件卸载动画 & 无法实现 & 可手动实现 \\
\rowcolor{tableheader}
包大小 & ~38KB（gzip） & 0 & ~60KB \\
性能 & 优化良好（使用transform） & 原生性能 & 优秀（专业级） \\
\rowcolor{tableheader}
\bottomrule
\end{tabular}
\end{table}

Framer Motion适合本项目：消息气泡入场、Live2D面板展开收起、按钮选项弹出等动画都可用声明式API简洁表达。

\subsection{前端技术选型汇总表}

\begin{table}[H]
\centering
\caption{前端技术栈最终选型清单}
\label{tab:frontend-summary}
\begin{tabular}{L{3.5cm}L{4cm}L{4.5cm}}
\toprule
\textbf{技术领域} & \textbf{选定技术} & \textbf{备选方案} \\
\midrule
\rowcolor{tableheader}
UI框架 & React 18 + TypeScript & Vue 3 + TypeScript \\
构建工具 & Vite 5.x & Webpack 5 \\
\rowcolor{tableheader}
状态管理 & Zustand 4.x & Redux Toolkit \\
样式方案 & Tailwind CSS 3.x + CSS Variables & CSS Modules + SCSS \\
\rowcolor{tableheader}
UI组件库 & 自定义 + Radix UI Primitives & shadcn/ui \\
Markdown渲染 & react-markdown + shiki & react-markdown + prismjs \\
\rowcolor{tableheader}
Live2D & Cubism Web SDK 4.x & 无（官方唯一） \\
WebSocket客户端 & 原生 WebSocket + 自定义封装 & Socket.io-client \\
\rowcolor{tableheader}
PWA & vite-plugin-pwa & 手动配置Workbox \\
动画库 & Framer Motion 11.x & CSS Transition + GSAP \\
\rowcolor{tableheader}
字体 & Inter + PingFang SC + JetBrains Mono & Noto Sans SC \\
图标 & Lucide React & Heroicons \\
\bottomrule
\end{tabular}
\end{table}

% ============================================================
\section{后端技术栈详解}
% ============================================================

\subsection{Web框架选型：FastAPI}

\begin{table}[H]
\centering
\caption{Python Web框架对比}
\label{tab:web-framework}
\begin{tabular}{L{2.5cm}L{3.5cm}L{3.5cm}L{3.5cm}}
\toprule
\textbf{维度} & \textbf{FastAPI} & \textbf{Flask} & \textbf{Django} \\
\midrule
\rowcolor{tableheader}
异步原生支持 & 完美（基于Starlette） & 需扩展（Flask 2.0+） & 有限（Django 3.1+） \\
WebSocket支持 & 原生内置 & 需Flask-SocketIO & 需Django Channels \\
\rowcolor{tableheader}
自动API文档 & OpenAPI + Swagger UI & 需Flask-RESTX等 & 需DRF \\
类型系统 & Pydantic深度集成 & 无原生 & 无原生 \\
\rowcolor{tableheader}
性能 & 极高（asyncio原生） & 中等 & 中等（重量级） \\
学习曲线 & 中等（需理解async） & 低 & 高（概念多） \\
\rowcolor{tableheader}
依赖注入 & 内置强大 & 无 & 无 \\
\bottomrule
\end{tabular}
\end{table}

\textbf{选择FastAPI的核心理由：}

\begin{enumerate}[leftmargin=2em]
\item \textbf{WebSocket原生支持}：FastAPI的\texttt{WebSocket}类可以直接在async函数中处理连接、发送和接收消息，无需额外依赖。
\item \textbf{Pydantic集成}：请求/响应的自动验证和序列化，与配置管理（Pydantic Settings）天然统一。
\item \textbf{异步架构}：管理Kimi CLI子进程（使用\texttt{asyncio.subprocess}）的同时，处理WebSocket连接，不会阻塞。
\item \textbf{自动生成文档}：局域网调试时，开发人员可以直接访问\texttt{/docs}查看和测试API。
\end{enumerate}

\subsection{ASGI服务器：Uvicorn}

\begin{table}[H]
\centering
\caption{ASGI服务器对比}
\label{tab:asgi-server}
\begin{tabular}{L{2.5cm}L{3.5cm}L{3.5cm}L{3.5cm}}
\toprule
\textbf{维度} & \textbf{Uvicorn} & \textbf{Daphne} & \textbf{Hypercorn} \\
\midrule
\rowcolor{tableheader}
HTTP性能 & 极高（基于uvloop + httptools） & 中等 & 高 \\
WebSocket性能 & 高 & 中等 & 高 \\
\rowcolor{tableheader}
HTTP/2支持 & 否（需gunicorn） & 否 & 是 \\
HTTP/3支持 & 否 & 否 & 实验性 \\
\rowcolor{tableheader}
进程管理 & 内置单进程，需gunicorn多进程 & 需外部管理 & 支持多进程 \\
兼容性 & 与FastAPI最佳 & Django Channels绑定 & 与Quart绑定深 \\
\rowcolor{tableheader}
\bottomrule
\end{tabular}
\end{table}

Uvicorn是FastAPI推荐的服务器，在局域网单用户场景下单进程完全足够。对于公网部署，可通过\texttt{gunicorn} + \texttt{uvicorn.workers.UvicornWorker}实现多进程。

\subsection{CLI桥接技术}

Kimi CLI的集成是本项目的核心挑战之一。CLI输出采用流式JSON格式，需要通过\texttt{asyncio.subprocess}进行非阻塞读取和解析。

\begin{lstlisting}[language=Python, caption=CLI桥接核心逻辑]
import asyncio
import json
import re
from typing import AsyncIterator

class KimiCLIBridge:
    def __init__(self, working_dir: str):
        self.working_dir = working_dir
        self.process: asyncio.subprocess.Process | None = None

    async def start(self):
        self.process = await asyncio.create_subprocess_exec(
            'kimi',
            'chat',
            cwd=self.working_dir,
            stdin=asyncio.subprocess.PIPE,
            stdout=asyncio.subprocess.PIPE,
            stderr=asyncio.subprocess.STDOUT,
        )

    async def send_message(self, text: str) -> AsyncIterator[dict]:
        if not self.process or self.process.stdin.is_closing():
            raise RuntimeError("CLI not started")

        self.process.stdin.write(f"{text}\n".encode())
        await self.process.stdin.drain()

        # 流式解析JSON输出
        while True:
            line = await self.process.stdout.readline()
            if not line:
                break
            yield self._parse_line(line.decode())

    def _parse_line(self, raw: str) -> dict:
        # 状态提示匹配
        if raw.startswith('◇'):
            return {"type": "status", "content": raw.strip()}
        # 选项列表
        if re.match(r'^\d+\.', raw):
            return {"type": "option", "content": raw.strip()}
        # 普通文本
        return {"type": "text", "content": raw}

    async def interrupt(self):
        """发送SIGINT打断当前生成"""
        if self.process:
            self.process.send_signal(2)  # SIGINT
\end{lstlisting}

\textbf{关键技术点：}
\begin{itemize}[leftmargin=2em]
\item \textbf{asyncio.subprocess}：\texttt{create\_subprocess\_exec}创建子进程，通过\texttt{PIPE}进行stdin/stdout通信，完全异步不阻塞事件循环。
\item \textbf{流式解析}：逐行读取stdout，通过正则/状态机识别不同输出类型（文本、状态、选项、JSON）。
\item \textbf{SIGINT处理}：\texttt{process.send\_signal(2)} 发送中断信号，Kimi CLI支持在生成过程中响应SIGINT并返回部分结果。
\item \textbf{异常处理}：需处理CLI崩溃、卡死、输出格式异常等边界情况。
\end{itemize}

\subsection{数据持久化：SQLite}

\begin{table}[H]
\centering
\caption{数据持久化方案对比}
\label{tab:database}
\begin{tabular}{L{2.5cm}L{3.5cm}L{3.5cm}L{3.5cm}}
\toprule
\textbf{维度} & \textbf{SQLite} & \textbf{PostgreSQL} & \textbf{MongoDB} \\
\midrule
\rowcolor{tableheader}
部署复杂度 & 零（文件数据库） & 需安装/配置服务 & 需安装服务 \\
配置要求 & 无 & 需调优 & 需配置 \\
\rowcolor{tableheader}
并发写入 & 单写多读（WAL模式） & 高并发优秀 & 高并发优秀 \\
数据规模 & GB级别适合 & TB级别适合 & TB级别适合 \\
\rowcolor{tableheader}
异步支持 & aiosqlite可用 & asyncpg优秀 & motor可用 \\
关系查询 & 支持SQL & 支持SQL & 不支持JOIN \\
\rowcolor{tableheader}
适用场景 & 本地单用户 & 多用户/服务器 & 非结构化数据 \\
\bottomrule
\end{tabular}
\end{table}

SQLite适用于阶段1（局域网单用户）：
\begin{itemize}[leftmargin=2em]
\item 无需安装，单文件存储，适合随应用一起分发。
\item \texttt{aiosqlite}提供异步接口，不阻塞FastAPI的事件循环。
\item WAL模式（Write-Ahead Logging）提升读写并发能力。
\item 阶段2（公网多用户）可平滑迁移到PostgreSQL，使用SQLAlchemy Async抽象数据库层。
\end{itemize}

\subsection{配置与数据验证：Pydantic Settings}

\begin{lstlisting}[language=Python, caption=Pydantic Settings配置示例]
from pydantic_settings import BaseSettings
from pydantic import Field

class ELAWSettings(BaseSettings):
    # 服务器配置
    host: str = Field(default="0.0.0.0", env="ELAW_HOST")
    port: int = Field(default=8765, env="ELAW_PORT")

    # CLI配置
    kimi_cli_path: str = Field(default="kimi", env="ELAW_KIMI_PATH")
    working_dir: str = Field(default=".", env="ELAW_WORK_DIR")

    # TTS配置
    tts_provider: str = Field(default="edge-tts", env="ELAW_TTS_PROVIDER")
    tts_voice: str = Field(default="zh-CN-XiaoxiaoNeural", env="ELAW_TTS_VOICE")

    # 情绪分析
    emotion_method: str = Field(default="keyword", env="ELAW_EMOTION_METHOD")

    class Config:
        env_file = ".env"
        env_file_encoding = "utf-8"
\end{lstlisting}

Pydantic Settings的优势：自动从环境变量和\texttt{.env}文件读取配置，类型验证在启动时完成，提供默认值降低首次配置门槛。

\subsection{情绪分析技术}

\begin{table}[H]
\centering
\caption{情绪分析方案对比}
\label{tab:emotion-analysis}
\begin{tabular}{L{2.5cm}L{3.5cm}L{3.5cm}L{3.5cm}}
\toprule
\textbf{维度} & \textbf{LLM-based} & \textbf{Embedding-based} & \textbf{Keyword-based} \\
\midrule
\rowcolor{tableheader}
准确率 & 最高（理解上下文） & 高（语义相似度） & 中等（字面匹配） \\
延迟 & 高（API调用~500ms） & 中（本地~100ms） & 低（本地~1ms） \\
\rowcolor{tableheader}
成本 & 高（API计费） & 低（一次性或本地） & 零（完全本地） \\
离线可用 & 否（需网络） & 是（本地模型） & 是（完全本地） \\
\rowcolor{tableheader}
实现复杂度 & 低（调用API） & 中（需管理向量） & 低（Trie树） \\
可扩展性 & 高（自动适应新表达） & 中（需扩充语料） & 低（需人工维护词库） \\
\rowcolor{tableheader}
\bottomrule
\end{tabular}
\end{table}

\textbf{推荐策略：三层级联}
\begin{enumerate}[leftmargin=2em]
\item \textbf{Keyword-based（首选）}：Trie树快速匹配，零延迟零成本，覆盖80\%常见场景。旧项目已有词库可复用。
\item \textbf{Embedding-based（补充）}：对未命中关键词的文本，计算与预定义情绪标签的向量相似度。可使用\texttt{text-embedding-3-small}或本地\texttt{sentence-transformers}。
\item \textbf{LLM-based（高质量场景）}：对特别重要或复杂的对话，调用Moonshot API进行精细情绪分析。
\end{enumerate}

\subsection{TTS技术}

\begin{table}[H]
\centering
\caption{TTS方案对比}
\label{tab:tts}
\begin{tabular}{L{2.5cm}L{3.5cm}L{3.5cm}L{3.5cm}}
\toprule
\textbf{维度} & \textbf{Edge TTS} & \textbf{百度/阿里API} & \textbf{本地TTS (Coqui)} \\
\midrule
\rowcolor{tableheader}
成本 & 免费 & 按量计费（~1元/千次） & 免费（GPU计算成本） \\
音质 & 高（Microsoft Edge语音） & 高（多种音色） & 中（取决于模型） \\
\rowcolor{tableheader}
延迟 & 中等（网络传输） & 中等（API调用） & 低（本地推理） \\
离线可用 & 否（需连接Microsoft服务） & 否 & 是 \\
\rowcolor{tableheader}
中文支持 & 优秀（XiaoxiaoNeural等） & 优秀 & 一般（模型依赖） \\
并发能力 & 高（无限制） & 受API配额限制 & 受GPU限制 \\
\rowcolor{tableheader}
\bottomrule
\end{tabular}
\end{table}

\textbf{选型：Edge TTS作为默认方案。} 它是Microsoft Edge浏览器的在线TTS服务，通过Python库\texttt{edge-tts}调用。音质优秀、免费、中文支持好，非常适合个人局域网使用。唯一的依赖是首次使用需要网络连接到Microsoft服务器。

\begin{lstlisting}[language=Python, caption=Edge TTS集成示例]
import edge_tts
import asyncio

async def generate_tts(text: str, output_file: str) -> str:
    communicate = edge_tts.Communicate(
        text=text,
        voice="zh-CN-XiaoxiaoNeural",
        rate="+0%",      # 语速
        volume="+0%",    # 音量
        pitch="+0Hz"     # 音调
    )
    await communicate.save(output_file)
    return output_file

# 获取音频数据用于Web Audio API分析
async def get_audio_bytes(text: str) -> bytes:
    communicate = edge_tts.Communicate(text, "zh-CN-XiaoxiaoNeural")
    audio_bytes = bytearray()
    async for chunk in communicate.stream():
        if chunk["type"] == "audio":
            audio_bytes.extend(chunk["data"])
    return bytes(audio_bytes)
\end{lstlisting}

\subsection{WebSocket服务端设计}

FastAPI的WebSocket原生支持：

\begin{lstlisting}[language=Python, caption=FastAPI WebSocket端点]
from fastapi import FastAPI, WebSocket, WebSocketDisconnect
from fastapi.staticfiles import StaticFiles
import json

app = FastAPI()

@app.websocket("/ws")
async def websocket_endpoint(websocket: WebSocket):
    await websocket.accept()
    client_id = id(websocket)
    try:
        while True:
            raw = await websocket.receive_text()
            message = json.loads(raw)

            if message["type"] == "chat":
                async for chunk in kimi_bridge.stream(message["text"]):
                    await websocket.send_json({
                        "type": "stream_chunk",
                        "data": chunk
                    })
            elif message["type"] == "interrupt":
                await kimi_bridge.interrupt()
                await websocket.send_json({
                    "type": "interrupted"
                })
            elif message["type"] == "ping":
                await websocket.send_json({"type": "pong"})
    except WebSocketDisconnect:
        print(f"Client {client_id} disconnected")
\end{lstlisting}

\textbf{并发模型说明：}
\begin{itemize}[leftmargin=2em]
\item FastAPI基于\texttt{asyncio}，每个WebSocket连接是一个独立的协程，单线程处理数万并发连接。
\item 在单用户场景下，仅有一个活跃WebSocket连接，资源消耗极低。
\item Kimi CLI子进程使用独立进程（非线程），通过asyncio PIPE通信，与WebSocket协程互不阻塞。
\end{itemize}

\subsection{安全与沙箱}

阶段1（局域网）的安全策略简化：
\begin{itemize}[leftmargin=2em]
\item 无用户认证（单用户/局域网信任）。
\item 保留AST审计（旧项目复用）：对Kimi CLI生成的代码进行抽象语法树分析，检测危险操作（如\texttt{rm -rf}、文件删除等）。
\item 保留敏感词过滤（旧项目复用）：对输入和输出进行Trie树过滤。
\item 文件上传限制：单文件10MB，限制扩展名（.jpg/.png/.txt/.pdf等），禁止可执行文件。
\item WebSocket无CORS限制（局域网内访问），阶段2公网部署时启用HTTPS/WSS和Token认证。
\end{itemize}

\subsection{后端技术选型汇总表}

\begin{table}[H]
\centering
\caption{后端技术栈最终选型清单}
\label{tab:backend-summary}
\begin{tabular}{L{3.5cm}L{4cm}L{4.5cm}}
\toprule
\textbf{技术领域} & \textbf{选定技术} & \textbf{备选方案} \\
\midrule
\rowcolor{tableheader}
Web框架 & FastAPI 0.110+ & Flask 2.x + Flask-SocketIO \\
ASGI服务器 & Uvicorn 0.27+ & Hypercorn \\
\rowcolor{tableheader}
进程管理 & asyncio.subprocess & threading + subprocess \\
数据存储 & SQLite + aiosqlite & PostgreSQL + asyncpg \\
\rowcolor{tableheader}
配置管理 & Pydantic Settings 2.x & python-dotenv + dataclasses \\
CLI解析 & 正则 + 状态机 & 纯JSON解析（需CLI支持） \\
\rowcolor{tableheader}
情绪分析 & 本地Trie + Embedding + LLM API & 仅LLM API \\
TTS & edge-tts & Azure TTS API \\
\rowcolor{tableheader}
日志 & structlog & Python标准logging \\
代码质量 & Ruff + MyPy & Flake8 + Black + MyPy \\
\rowcolor{tableheader}
测试 & pytest + pytest-asyncio & unittest \\
\bottomrule
\end{tabular}
\end{table}

% ============================================================
\section{部署与运维架构}
% ============================================================

\subsection{局域网部署架构}

\begin{figure}[H]
\centering
\begin{tikzpicture}[
    node distance=1.5cm and 2.5cm,
    box/.style={rectangle, rounded corners, draw=primaryblue, fill=lightbg, minimum width=2.8cm, minimum height=1cm, align=center, font=\small},
    device/.style={rectangle, rounded corners, draw=gray, fill=white, minimum width=2.5cm, minimum height=0.8cm, align=center, font=\small},
    arrow/.style={-{Stealth[scale=0.8]}, thick, primaryblue},
    net/.style={rectangle, dashed, draw=gray, rounded corners, inner sep=0.4cm},
    label/.style={font=\tiny\color{gray}}
]

% 局域网
\node[net, fit={(0,0) (8,3.5)}, label={[label, anchor=south]above:局域网 (WiFi)}] (lan) {};

% 服务器
\node[box] (server) at (2, 2) {ELAW Server\\{\tiny Python + FastAPI}};
\node[box, below=0.5cm of server] (cli) {Kimi Code CLI\\{\tiny subprocess}};

% 客户端
\node[device, right=3cm of server] (desktop) {桌面浏览器\\{Chrome/Firefox}};
\node[device, below=0.5cm of desktop] (mobile) {手机浏览器\\{Safari/Chrome}};

% 连接
\draw[arrow] (desktop) -- node[above, font=\tiny] {WebSocket} (server);
\draw[arrow] (mobile) -- node[below, font=\tiny] {WebSocket} (server);
\draw[arrow, <->] (server) -- node[right, font=\tiny] {PIPE} (cli);

% 端口标注
\node[label, below=0.1cm of server] {端口 8765};
\node[label, right=0.1cm of desktop] {PWA};
\node[label, right=0.1cm of mobile] {PWA};

\end{tikzpicture}
\caption{局域网部署架构图}
\label{fig:lan-arch}
\end{figure}

\textbf{局域网部署要点：}
\begin{itemize}[leftmargin=2em]
\item \textbf{IP发现}：服务器启动时打印本机所有IP地址（\texttt{0.0.0.0:8765}），用户可直接在内网设备访问。
\item \textbf{端口配置}：默认8765，可通过环境变量\texttt{ELAW\_PORT}修改。避免使用80/443（需要root权限）。
\item \textbf{PWA安装}：手机浏览器访问后，Chrome（Android）会提示"添加到主屏幕"；iOS Safari需点击分享按钮手动添加。
\item \textbf{防火墙}：Windows/macOS首次运行需放行网络请求；确保手机和电脑在同一WiFi下。
\end{itemize}

\subsection{公网升级路径}

阶段2公网部署架构：

\begin{figure}[H]
\centering
\begin{tikzpicture}[
    node distance=1.5cm and 2.5cm,
    box/.style={rectangle, rounded corners, draw=primaryblue, fill=lightbg, minimum width=2.5cm, minimum height=0.8cm, align=center, font=\small},
    cloud/.style={ellipse, draw=gray, dashed, fill=white, minimum width=2cm, minimum height=1cm, align=center, font=\small},
    arrow/.style={-{Stealth[scale=0.8]}, thick, primaryblue}
]

\node[cloud] (cdn) at (0, 2) {CDN\\{静态资源}};
\node[cloud] (dns) at (0, 0) {DNS\\{域名解析}};
\node[box] (nginx) at (3, 1) {Nginx\\{反向代理}};
\node[box] (server) at (6, 1) {ELAW Server};
\node[box] (db) at (6, -0.5) {PostgreSQL};
\node[cloud, right=2cm of nginx] (client) at (9, 1) {客户端\\{浏览器}};

\draw[arrow] (client) -- node[above, font=\tiny] {HTTPS/WSS} (nginx);
\draw[arrow] (nginx) -- node[above, font=\tiny] {proxy} (server);
\draw[arrow] (server) -- node[right, font=\tiny] {SQL} (db);
\draw[arrow] (dns) -- (nginx);
\draw[arrow, dashed] (cdn) -- (client);

\end{tikzpicture}
\caption{公网部署架构图}
\label{fig:wan-arch}
\end{figure}

\textbf{公网升级步骤：}
\begin{enumerate}[leftmargin=2em]
\item \textbf{Docker化}：多阶段构建（Node镜像构建前端，Python镜像运行后端）。
\item \textbf{反向代理}：Nginx处理HTTPS终止、静态文件服务、WebSocket转发。
\item \textbf{SSL证书}：Let's Encrypt自动续期（\texttt{certbot}）。
\item \textbf{域名与备案}：国内服务器需ICP备案；海外服务器无需备案。
\item \textbf{数据库迁移}：SQLite导出后导入PostgreSQL，SQLAlchemy兼容两者。
\end{enumerate}

\subsection{移动端适配要点}

\begin{table}[H]
\centering
\caption{移动端适配问题与解决方案}
\label{tab:mobile-adapt}
\begin{tabular}{L{3cm}L{4.5cm}L{4.5cm}}
\toprule
\textbf{问题} & \textbf{现象} & \textbf{解决方案} \\
\midrule
\rowcolor{tableheader}
WebSocket断开 & 手机息屏或切换App后连接断开 & 实现重连机制 + 恢复会话状态 \\
PWA保活 & 后台运行被系统清理 & 前端保存未发送消息，重连后补发 \\
\rowcolor{tableheader}
软键盘遮挡 & 输入框被键盘覆盖 & 使用\texttt{visual viewport}监听调整布局 \\
WebGL性能 & 低端手机掉帧 & 提供画质选项（高/中/低/关闭） \\
\rowcolor{tableheader}
屏幕尺寸 & 手机小屏、桌面大屏 & 响应式布局 + 断点设计 \\
\bottomrule
\end{tabular}
\end{table}

% ============================================================
\section{开发工具链与工程规范}
% ============================================================

\subsection{代码质量工具}

\begin{table}[H]
\centering
\caption{代码质量工具配置}
\label{tab:code-quality}
\begin{tabular}{L{2.5cm}L{3cm}L{5cm}L{2.5cm}}
\toprule
\textbf{语言} & \textbf{工具} & \textbf{作用} & \textbf{配置方式} \\
\midrule
\rowcolor{tableheader}
Python & Ruff & Lint + Format（替代Flake8 + Black） & pyproject.toml \\
Python & MyPy & 静态类型检查 & pyproject.toml \\
\rowcolor{tableheader}
TypeScript & ESLint & 代码规范检查 & .eslintrc.json \\
TypeScript & Prettier & 代码格式化 & .prettierrc.json \\
\rowcolor{tableheader}
TypeScript & TypeScript & 编译时类型检查 & tsconfig.json \\
\bottomrule
\end{tabular}
\end{table}

\textbf{Ruff配置示例（pyproject.toml）：}
\begin{lstlisting}[language=Python, caption=Ruff配置]
[tool.ruff]
line-length = 88
target-version = "py311"

[tool.ruff.lint]
select = ["E", "F", "I", "N", "W", "UP", "B", "C4", "SIM"]
ignore = ["E501"]

[tool.ruff.lint.pydocstyle]
convention = "google"
\end{lstlisting}

\subsection{测试策略}

\begin{table}[H]
\centering
\caption{测试分层策略}
\label{tab:testing}
\begin{tabular}{L{2.5cm}L{3cm}L{4cm}L{3.5cm}}
\toprule
\textbf{层级} & \textbf{工具} & \textbf{覆盖范围} & \textbf{执行频率} \\
\midrule
\rowcolor{tableheader}
单元测试 & pytest (后端) & 工具函数、解析器、状态机 & 每次提交 \\
单元测试 & Vitest (前端) & 组件逻辑、Hooks、工具函数 & 每次提交 \\
\rowcolor{tableheader}
集成测试 & pytest-asyncio & WebSocket协议、CLI桥接 & 每次PR \\
E2E测试 & Playwright & 完整用户流程（PWA安装、聊天、打断） & 每日/发布前 \\
\rowcolor{tableheader}
\bottomrule
\end{tabular}
\end{table}

\subsection{版本管理与CI/CD}

\textbf{Git工作流：}采用GitHub Flow（简化版GitFlow），主分支为\texttt{main}，功能开发从\texttt{main}切出分支，通过PR合并。

\textbf{GitHub Actions 流水线：}
\begin{enumerate}[leftmargin=2em]
\item \textbf{Lint检查}：Ruff、ESLint、TypeScript编译检查。
\item \textbf{单元测试}：pytest + Vitest，要求覆盖率$>$80\%。
\item \textbf{构建验证}：Vite构建、Docker镜像构建。
\item \textbf{E2E测试}：Playwright在Docker中运行完整流程。
\end{enumerate}

\subsection{项目目录结构建议}

\begin{lstlisting}[caption=推荐目录结构]
elaw/
  backend/
    elaw_server/
      __init__.py
      main.py          # FastAPI入口
      config.py        # Pydantic Settings
      websocket.py     # WebSocket处理
      cli_bridge.py    # Kimi CLI桥接
      tts.py           # TTS服务
      emotion.py       # 情绪分析
      models.py        # 数据模型/SQLAlchemy
      security.py      # 安全/审计
      utils.py         # 工具函数
    tests/
    pyproject.toml
    requirements.txt
  frontend/
    src/
      main.tsx
      App.tsx
      components/      # UI组件
      hooks/           # 自定义Hooks
      stores/          # Zustand状态
      lib/             # 工具函数
      websocket/       # WebSocket客户端
      live2d/          # Live2D封装
      styles/          # 全局样式
    public/              # 静态资源、manifest
    tests/
    vite.config.ts
    tailwind.config.js
    package.json
  docker/
    Dockerfile.backend
    Dockerfile.frontend
    docker-compose.yml
  assets/                  # Live2D模型、表情包
  README.md
\end{lstlisting}

% ============================================================
\section{技术风险评估与缓解方案}
% ============================================================

\subsection{风险矩阵}

\begin{table}[H]
\centering
\caption{技术风险矩阵}
\label{tab:risk-matrix}
\begin{tabular}{L{3.5cm}C{1.8cm}C{1.8cm}C{1.8cm}L{4.5cm}}
\toprule
\textbf{风险项} & \textbf{影响} & \textbf{概率} & \textbf{综合} & \textbf{缓解措施} \\
\midrule
\rowcolor{tableheader}
WebGL兼容性（老旧设备） & 高 & 中 & \riskcell{2} & 提供WebGL检测，降级为静态图片 \\
WebSocket稳定性（移动端） & 高 & 高 & \riskcell{3} & 重连机制 + 会话状态持久化 \\
\rowcolor{tableheader}
TTS延迟（网络依赖） & 中 & 中 & \riskcell{2} & 预缓存常用回复、异步预生成 \\
Live2D模型版权 & 中 & 低 & \riskcell{1} & 使用官方免费模型或自购授权 \\
\rowcolor{tableheader}
CLI输出格式变化 & 高 & 中 & \riskcell{2} & 版本锁定、解析器容错设计 \\
Live2D WebGL性能（低端手机） & 中 & 高 & \riskcell{2} & 画质分级、模型LOD优化 \\
\rowcolor{tableheader}
PWA离线缓存失效 & 低 & 中 & \riskcell{1} & 版本号管理、强制刷新机制 \\
局域网IP变动 & 低 & 高 & \riskcell{1} & mDNS广播（可选）、启动时打印IP \\
\rowcolor{tableheader}
Safari PWA限制 & 中 & 中 & \riskcell{2} & 引导用户手动安装、提供Safari兼容性说明 \\
并发TTS请求拥塞 & 中 & 低 & \riskcell{1} & 队列管理、单音频流限制 \\
\bottomrule
\end{tabular}
\end{table}

\subsection{重点风险详细分析}

\subsubsection{WebSocket稳定性风险（最高优先级）}

移动端WebSocket连接极易因以下原因断开：
\begin{itemize}[leftmargin=2em]
\item 手机息屏（iOS约30秒关闭连接）
\item 网络切换（WiFi与蜂窝数据切换）
\item 应用切换（后台被系统清理）
\item 防火墙/NAT超时（某些路由器会断开空闲连接）
\end{itemize}

\textbf{缓解方案：}
\begin{enumerate}[leftmargin=2em]
\item 客户端实现指数退避重连（最多10次，最大间隔30秒）。
\item 心跳检测：每30秒ping-pong，防止NAT超时。
\item 会话状态恢复：后端维护会话状态，重连后发送完整会话历史和当前生成状态。
\item 消息队列：用户输入在断线时缓存，重连后自动发送。
\item 用户反馈：断线时显示"连接中..."状态，重连成功后恢复。
\end{enumerate}

\subsubsection{CLI输出格式不稳定风险}

Kimi Code CLI可能在新版本中改变输出格式（如新增字段、改变状态提示图标）。

\textbf{缓解方案：}
\begin{enumerate}[leftmargin=2em]
\item 解析器设计为容错模式：未知字段忽略，不阻断流程。
\item 版本锁定：文档中标注兼容的CLI版本范围。
\item 正则 + 状态机结合：对文本输出使用模式匹配，而非严格JSON Schema。
\item 灰度测试：CLI更新后先本地验证再发布。
\item 提供降级方案：解析失败时直接显示原始文本，不中断服务。
\end{enumerate}

% ============================================================
\section{总结与推荐配置}
% ============================================================

\subsection{最小可行产品 (MVP) 技术栈}

MVP目标：最快两周内可运行的原型，覆盖核心聊天、Live2D展示、基础TTS。

\begin{table}[H]
\centering
\caption{MVP推荐技术栈清单}
\label{tab:mvp-stack}
\begin{tabular}{L{3.5cm}L{4.5cm}L{4cm}}
\toprule
\textbf{技术领域} & \textbf{推荐技术} & \textbf{版本} \\
\midrule
\rowcolor{tableheader}
前端框架 & React 18 + TypeScript & >=18.2.0 \\
构建工具 & Vite & >=5.0.0 \\
\rowcolor{tableheader}
状态管理 & Zustand & >=4.5.0 \\
样式 & Tailwind CSS & >=3.4.0 \\
\rowcolor{tableheader}
Markdown & react-markdown + shiki & >=9.0 / >=1.0 \\
WebSocket & 原生 WebSocket + 封装 & 内置 \\
\rowcolor{tableheader}
Live2D & Cubism Web SDK 4.x & SDK 4.2+ \\
动画 & Framer Motion & >=11.0 \\
\rowcolor{tableheader}
后端框架 & FastAPI & >=0.110 \\
ASGI服务器 & Uvicorn & >=0.27 \\
\rowcolor{tableheader}
数据存储 & SQLite (aiosqlite) & >=3.12 \\
TTS & edge-tts & >=6.1 \\
\rowcolor{tableheader}
配置管理 & Pydantic Settings & >=2.2 \\
情绪分析 & 本地Trie（关键词） & 旧项目复用 \\
\rowcolor{tableheader}
代码质量 & Ruff + MyPy & 最新 \\
测试 & pytest + Vitest & 最新 \\
\bottomrule
\end{tabular}
\end{table}

\subsection{完整产品技术栈}

完整目标：包含所有高级功能（PWA、Embedding情绪分析、多音色TTS、E2E测试、公网部署）。

\begin{table}[H]
\centering
\caption{完整产品推荐技术栈清单}
\label{tab:full-stack}
\begin{tabular}{L{3.5cm}L{4.5cm}L{4cm}}
\toprule
\textbf{技术领域} & \textbf{推荐技术} & \textbf{版本/备注} \\
\midrule
\rowcolor{tableheader}
前端框架 & React 18 + TypeScript & >=18.2.0 \\
PWA & vite-plugin-pwa + Workbox & >=0.19.0 \\
\rowcolor{tableheader}
Live2D & Cubism Web SDK 5.x（如有） & 跟进最新版本 \\
图表渲染 & Mermaid.js + react-markdown & 支持工具结果展示 \\
\rowcolor{tableheader}
Embedding情绪 & sentence-transformers (本地) & 支持离线语义分析 \\
LLM情绪 & Moonshot API (moonshot-v1-8k) & 高质量场景 \\
\rowcolor{tableheader}
TTS备选 & 百度/阿里 TTS API & 多音色选择 \\
STT（未来） & Whisper (本地) & OpenAI开源 \\
\rowcolor{tableheader}
公网数据库 & PostgreSQL + asyncpg & 多用户场景 \\
容器化 & Docker + Docker Compose & 多阶段构建 \\
\rowcolor{tableheader}
反向代理 & Nginx + Let's Encrypt & HTTPS/WSS \\
CI/CD & GitHub Actions & 完整流水线 \\
\rowcolor{tableheader}
E2E测试 & Playwright & PWA流程测试 \\
\bottomrule
\end{tabular}
\end{table}

\subsection{技术栈演进路线图}

\begin{figure}[H]
\centering
\begin{tikzpicture}[
    stage/.style={rectangle, rounded corners, draw=primaryblue, fill=lightbg, minimum width=3.5cm, minimum height=1.2cm, align=center, font=\small},
    arrow/.style={-{Stealth[scale=0.8]}, thick, primaryblue},
    timeline/.style={rectangle, draw=gray, dashed, rounded corners, inner sep=0.3cm}
]

\node[timeline, fit={(-0.5,-0.5) (12,1.5)}] (t1) {};
\node[label, anchor=south] at (t1.north) {阶段1：MVP（2-4周）};
\node[stage] at (1.5, 0.5) {React + Vite\\FastAPI + SQLite};
\node[stage] at (5.5, 0.5) {WebSocket基础\\Live2D展示};
\node[stage] at (9.5, 0.5) {Edge TTS\\关键词情绪};

\node[timeline, fit={(-0.5,-2.5) (12,-0.5)}] (t2) {};
\node[label, anchor=south] at (t2.north) {阶段2：完整版（4-8周）};
\node[stage] at (1.5, -1.5) {PWA + Framer Motion\\响应式优化};
\node[stage] at (5.5, -1.5) {Embedding情绪\\shiki高亮};
\node[stage] at (9.5, -1.5) {打断机制\\会话持久化};

\node[timeline, fit={(-0.5,-4.5) (12,-2.5)}] (t3) {};
\node[label, anchor=south] at (t3.north) {阶段3：公网部署（可选）};
\node[stage] at (2.5, -3.5) {Docker化\\Nginx反向代理};
\node[stage] at (6.5, -3.5) {PostgreSQL\\用户认证};
\node[stage] at (10.5, -3.5) {HTTPS/WSS\\CDN加速};

\draw[arrow] (t1) -- (t2);
\draw[arrow] (t2) -- (t3);

\end{tikzpicture}
\caption{技术栈演进路线图}
\label{fig:roadmap}
\end{figure}

\vspace{1cm}
\noindent\textbf{结语：}ELAW项目的技术栈设计以"本地化优先、渐进增强"为原则。阶段1（MVP）追求最小依赖和最快开发速度，通过React + Vite + FastAPI + SQLite的组合，可在2-4周内完成核心功能。随着需求演进，逐步引入PWA、Embedding分析、Docker化等高级能力，最终形成可公网部署的完整产品。所有技术选型都经过与旧项目（ACP-QQ-Bridge）的经验对比验证，确保技术路径的可行性和可维护性。

\end{document}
